\hsection{Chapter 3}
\sol{3.1.3} Only (a) has the form $\forall x\in X,\ \exists!\,y\in Y,\ p(x,y)$: ``for each real number $a$, the equation $x^2+2ax+a^2=0$ has exactly one real solution $x$'' is $\forall a\in\R,\ \exists!\,x\in\R,\ x^2+2ax+a^2=0$. Since $x^2+2ax+a^2=(x+a)^2$, the unique solution is $x=-a$, so the function is $\R\to\R$, $a\mapsto-a$. Statement (b) begins with $\exists!\,a\in\R$ and (c) with $\exists!\,n\in\N$, so neither has the required form.

\sol{3.1.22} By Definition 3.1.17, $g\circ f$ requires the codomain of $f$ to equal the domain of $g$, and $Y\ne Z$. Let $i:Y\to Z$ be the inclusion, $i(y)=y$ (well-defined since $Y\subseteq Z$). Then $h=g\circ i\circ f=(g\circ i)\circ f$: for $x\in X$, $(g\circ(i\circ f))(x)=g(i(f(x)))=g(f(x))=h(x)$, and the domains and codomains agree ($X\to W$).

\sol{3.1.28} Let $x\in X$. \emph{(b)} If $x\in A$ and $x\in B$: right side $1+1-1=1$, and $x\in A\cup B$ so left side $1$. If $x\in A$, $x\notin B$: $1+0-0=1$ and $x\in A\cup B$. If $x\notin A$, $x\in B$: symmetric. If $x\notin A$, $x\notin B$: $0+0-0=0$ and $x\notin A\cup B$. In all cases $\chi_{A\cup B}(x)=\chi_A(x)+\chi_B(x)-\chi_A(x)\chi_B(x)$. \emph{(c)} If $x\in A$ then $x\notin X\setminus A$, so $\chi_{X\setminus A}(x)=0=1-1$; if $x\notin A$ then $x\in X\setminus A$ (as $x\in X$), so $\chi_{X\setminus A}(x)=1=1-0$.

\sol{3.1.30} For $x\in X$, by Theorem 3.1.27: $\chi_{X\setminus(A\cup B)}=1-\chi_{A\cup B}=1-\chi_A-\chi_B+\chi_A\chi_B=(1-\chi_A)(1-\chi_B)=\chi_{X\setminus A}\chi_{X\setminus B}=\chi_{(X\setminus A)\cap(X\setminus B)}$, so $X\setminus(A\cup B)=(X\setminus A)\cap(X\setminus B)$ by Strategy 3.1.26. Likewise $\chi_{(X\setminus A)\cup(X\setminus B)}=(1-\chi_A)+(1-\chi_B)-(1-\chi_A)(1-\chi_B)=2-\chi_A-\chi_B-(1-\chi_A-\chi_B+\chi_A\chi_B)=1-\chi_A\chi_B=1-\chi_{A\cap B}=\chi_{X\setminus(A\cap B)}$, giving $X\setminus(A\cap B)=(X\setminus A)\cup(X\setminus B)$.

\sol{3.1.34} If $y\in f[\varnothing]$ then $y=f(x)$ for some $x\in\varnothing$, which is impossible. So $f[\varnothing]$ has no elements: $f[\varnothing]=\varnothing$.

\sol{3.1.42} Let $A=f^{-1}[\{1\}]=\{x\in X\mid f(x)=1\}$. Both $f$ and $\chi_A$ have domain $X$ and codomain $\{0,1\}=[2]$. Let $x\in X$. If $f(x)=1$ then $x\in A$, so $\chi_A(x)=1=f(x)$. If $f(x)\ne1$ then $f(x)=0$ (the only other element of the codomain) and $x\notin A$, so $\chi_A(x)=0=f(x)$. Hence $f=\chi_A$ by function extensionality.

\sol{3.2.7} $f:\N\to\Z$, $f(n)=n^2$: injective. If $m,n\in\N$ and $m^2=n^2$ then $(m-n)(m+n)=0$; if $m+n=0$ then $m=n=0$; otherwise $m-n=0$. Either way $m=n$. \ $g:\Z\to\N$, $g(n)=n^2$: not injective, $g(-1)=1=g(1)$. \ $h(x,y,z)=2^x3^y5^z$: injective, by uniqueness of prime factorisation (the fundamental theorem of arithmetic, Theorem 7.2.12, which the book also uses in Example 3.2.24): if $2^x3^y5^z=2^{x'}3^{y'}5^{z'}$ then the exponents of each prime agree, so $(x,y,z)=(x',y',z')$.

\sol{3.2.13} Take $V=f[X]$ and $g:X\to f[X]$, $g(x)=f(x)$ (well-defined since $f(x)\in f[X]$). For $y\in f[X]$ there is $x\in X$ with $f(x)=y$ (Definition 3.1.31), so $g(x)=y$: $g$ is surjective.

\sol{3.2.15} Take $Z=f[X]$, $p:X\to f[X]$, $p(x)=f(x)$ (surjective by Exercise 3.2.13) and $i:f[X]\to Y$, $i(y)=y$ (injective: $i(y)=i(y')$ means $y=y'$). Then $(i\circ p)(x)=i(f(x))=f(x)$ for all $x\in X$, so $f=i\circ p$.

\sol{3.2.21} ($\Rightarrow$) Suppose $f$ injective and let $y\in f[X]$. Existence: $y=f(x)$ for some $x\in X$ by definition of $f[X]$. Uniqueness: if $y=f(x)=f(x')$ then $x=x'$ by injectivity. ($\Leftarrow$) Suppose the condition holds and let $a,b\in X$ with $f(a)=f(b)=:y$. Then $y\in f[X]$, and both $a$ and $b$ are elements $x$ with $y=f(x)$; by uniqueness $a=b$.

\sol{3.2.22} Injective: if $2^m3^n=2^{m'}3^{n'}$ then $m=m'$ and $n=n'$ by uniqueness of prime factorisation (Theorem 7.2.12). (Directly: if $m<m'$, dividing by $2^m$ gives $3^n=2^{m'-m}3^{n'}$ with left side odd and right side even, impossible; so $m=m'$, then $3^n=3^{n'}$ gives $n=n'$.) Decoding: $1=2^03^0\mapsto(0,0)$; $24=2^3\cdot3\mapsto(3,1)$; $7776=6^5=2^53^5\mapsto(5,5)$; $59049=3^{10}\mapsto(0,10)$; $396718580736=2^{10}\cdot3^{18}\mapsto(10,18)$.

\sol{3.2.28} Let $x\in X$. Then $x=g(f(x))$, so $x$ is the value of $g$ at $f(x)\in Y$. Hence $g$ is surjective.

\sol{3.2.31} Let $g:Y\to X$ with $f\circ g=\id_Y$ and let $y\in Y$. Then $x=g(y)\in X$ satisfies $f(x)=f(g(y))=y$. So $f$ is surjective.

\sol{3.2.39(c)} Every $k\in\N$ has $k+1\ge1$, and every positive integer can be written uniquely as $2^m q$ with $m\in\N$ and $q$ odd (fundamental theorem of arithmetic: $m$ is the exponent of $2$ and $q$ the product of the odd prime powers; directly: if $2^mq=2^{m'}q'$ with $q,q'$ odd and $m<m'$ then $q=2^{m'-m}q'$ is even, contradiction, so $m=m'$ and $q=q'$). Write $q=2n+1$ and define $h^{-1}:\N\to\N\times\N$ by $h^{-1}(k)=(m,n)$ where $k+1=2^m(2n+1)$. Then $h^{-1}(h(m,n))=h^{-1}(2^m(2n+1)-1)=(m,n)$ by uniqueness, and $h(h^{-1}(k))=2^m(2n+1)-1=(k+1)-1=k$. So $h^{-1}$ is a two-sided inverse.

\sol{3.E.1} (a) Yes: $x+y=1\Leftrightarrow y=1-x$; $f(x)=1-x$. (b) No: for $x=0$ both $y=1$ and $y=-1$ satisfy the equation, so no single value $f(0)$ works (and for $x=2$ no $y$ at all). (c) No: for $x=0$ the equation ``$x=0$'' is true for every $y$, but $y=f(0)$ holds for only one $y$. (d) Yes: $f(x)=0$ for all $x$. (e) Yes: $1+x^2\ne0$, so $(1+x^2)y=1\Leftrightarrow y=\frac1{1+x^2}$. (f) No: for $x=1$ every $y$ satisfies $(1-1)y=0$.

\sol{3.E.3} Let $f:X\to\varnothing$ be a function. If $X$ had an element $a$, then $f(a)\in\varnothing$, impossible; so $X=\varnothing$. Any two functions $\varnothing\to\varnothing$ have the same domain and codomain and (vacuously) the same values, so they are equal by Axiom 3.1.4. Hence there is exactly one function with empty codomain, the empty function $\varnothing\to\varnothing$; its domain is $\varnothing$.

\sol{3.E.5} $f(-x)=(-x)^n=(-1)^nx^n$. If $n$ is even, $(-1)^n=1$ so $f(-x)=f(x)$: $f$ is even; if $n$ is odd, $(-1)^n=-1$ so $f(-x)=-f(x)$: $f$ is odd. Conversely, if $f$ is even then $f(-1)=f(1)$ gives $(-1)^n=1$, so $n$ is even (for odd $n$, $(-1)^n=-1\ne1$); if $f$ is odd then $(-1)^n=-1$, so $n$ is odd.

\sol{3.E.6} If $f$ is even and odd then for all $x$, $f(x)=f(-x)=-f(x)$, so $2f(x)=0$, $f(x)=0$. Thus $f$ is the zero function, which is indeed even and odd ($0=0=-0$). So exactly one such function exists.

\sol{3.E.7} Write $-U=\{-u\mid u\in U\}$; note $x\in-U\Leftrightarrow-x\in U$ (if $x=-u$ with $u\in U$ then $-x=u\in U$; conversely $x=-(-x)$). (a) $\chi_U$ even $\Leftrightarrow\forall x,\ \chi_U(-x)=\chi_U(x)\Leftrightarrow\forall x,\ (-x\in U\Leftrightarrow x\in U)\Leftrightarrow\forall x,\ (x\in-U\Leftrightarrow x\in U)\Leftrightarrow-U=U$. (b) $\chi_U$ odd $\Leftrightarrow\forall x,\ \chi_U(-x)=-\chi_U(x)\Leftrightarrow\forall x,\ (-x\in U\Leftrightarrow x\notin U)\Leftrightarrow\forall x,\ (x\in-U\Leftrightarrow x\in\R\setminus U)\Leftrightarrow-U=\R\setminus U$.

\sol{3.E.8} Put $g(x)=\frac{f(x)+f(-x)}2$ and $h(x)=\frac{f(x)-f(-x)}2$. Then $g(-x)=g(x)$, $h(-x)=-h(x)$ and $g(x)+h(x)=f(x)$. Uniqueness: if also $f=g'+h'$ with $g'$ even, $h'$ odd, then $g-g'=h'-h$; the left side is even and the right side is odd, so this function is both, hence zero by Exercise 3.E.6. So $g=g'$ and $h=h'$.

\sol{3.E.9} For $n=0$: $[0]=\varnothing$, and $\theta_0:\varnothing\to\varnothing$ is the empty function, which is $\id_\varnothing$. For $n\ge1$: first, $\theta_1=\id_{[1]}$, since $[1]=\{0\}$ has only one function to itself (3.E.39). Fix $k\in[n]$ and let $f:[1]\to[n]$, $f(0)=k$. Then $f\circ\theta_1=\theta_n\circ f$ gives $k=f(0)=f(\theta_1(0))=\theta_n(f(0))=\theta_n(k)$. As $k$ was arbitrary, $\theta_n=\id_{[n]}$.

\sol{3.E.10} $\chi_{U\symdiff V}=\chi_U+\chi_V-2\chi_U\chi_V$. Check by cases: $x$ in both: $1+1-2=0$ and $x\notin U\symdiff V$; in exactly one: $1+0-0=1$ and $x\in U\symdiff V$; in neither: $0$. (Equivalently $(\chi_U-\chi_V)^2$.)

\sol{3.E.11} On $\R\setminus\{0\}=(-\infty,0)\cup(0,\infty)$, a function $F$ with $F'(x)=\frac1x$ satisfies $(F(x)-\log|x|)'=0$ separately on each of the two intervals, so $F-\log|x|$ is constant on $(-\infty,0)$ (say $a$) and constant on $(0,\infty)$ (say $b$), and these constants are independent. Hence $F(x)=\log|x|+a\chi_U(x)+b\chi_V(x)$ with $U=(-\infty,0)$, $V=(0,\infty)$; conversely every such $F$ has derivative $\frac1x$.

\sol{3.E.12} $f[\R]=[1,\infty)$. ($\subseteq$) $1+x^2\ge1$ so $\sqrt{1+x^2}\ge1$. ($\supseteq$) For $y\ge1$ put $x=\sqrt{y^2-1}\in\R$; then $f(x)=\sqrt{1+y^2-1}=y$.

\sol{3.E.14} $f(1)=\{0\}$, $f(2)=\{0,1\}$, and in general $f(n)=[n]=\{0,1,\dots,n-1\}$ (induction: $f(n+1)=[n]\cup\{n\}=[n+1]$). So $f[\N]=\{[n]\mid n\in\N\}=\{\varnothing,\{0\},\{0,1\},\{0,1,2\},\dots\}$.

\sol{3.E.15} $f[\R^\R]$ is the set of even functions $\R\to\R$. ($\subseteq$) $f(h)(-x)=h(|-x|)=h(|x|)=f(h)(x)$, so every $f(h)$ is even. ($\supseteq$) Let $g$ be even. Then $f(g)(x)=g(|x|)=g(x)$ for all $x$ (if $x\ge0$ then $|x|=x$; if $x<0$ then $g(|x|)=g(-x)=g(x)$), so $f(g)=g$ and $g\in f[\R^\R]$.

\sol{3.E.16} $\sqrt{1+x^2}\in(-5,5]$ iff $\sqrt{1+x^2}\le5$ (it is always $>0$) iff $1+x^2\le25$ iff $x^2\le24$ iff $-2\sqrt6\le x\le2\sqrt6$. So $f^{-1}[(-5,5]]=[-2\sqrt6,2\sqrt6]$.

\sol{3.E.18} $f^{-1}[\{\varnothing\}]=\{(A,B)\in\PS(\N)\times\PS(\N)\mid A\cap B=\varnothing\}$, the set of pairs of disjoint subsets of $\N$.

\sol{3.E.20} \emph{Reading:} $v:A\to f[X]$ should satisfy $v(a)=i(a)$ for all $a\in A$. Let $q:X\to f[X]$, $q(x)=f(x)$, and $j:f[X]\to Y$, $j(y)=y$; $j$ is injective and $j\circ q=f$. By (iii) there is a unique $u:A\to f[X]$ with $j\circ u=i$, i.e.\ $u(a)=i(a)$ for all $a\in A$. So $v=u$ is the only candidate (uniqueness). $v$ is injective: $v(a)=v(b)\Rightarrow i(a)=i(b)\Rightarrow a=b$ by (i). $v$ is surjective: let $y\in f[X]$, say $y=f(x)$; by (ii) $y=i(p(x))=v(p(x))$. So $v$ is the unique bijection with $v(a)=i(a)$.

\sol{3.E.21} This is Exercise 3.1.40 (worked in Section A).

\sol{3.E.22} (a) $z\in(g\circ f)[U]\Leftrightarrow\exists x\in U,\ z=g(f(x))$; and $z\in g[f[U]]\Leftrightarrow\exists y\in f[U],\ z=g(y)\Leftrightarrow\exists x\in U,\ z=g(f(x))$ (unfold $y\in f[U]$). (b) $x\in(g\circ f)^{-1}[W]\Leftrightarrow g(f(x))\in W\Leftrightarrow f(x)\in g^{-1}[W]\Leftrightarrow x\in f^{-1}[g^{-1}[W]]$.

\sol{3.E.23} ($\subseteq$) Let $(x,z)\in p[S]$; then $(x,z)=p(x,y,z)$ for some $(x,y,z)\in S$, so $y=f(x)$ and $z=g(y)=g(f(x))=(g\circ f)(x)$; hence $(x,z)\in\mathrm{Gr}(g\circ f)$. ($\supseteq$) Let $(x,(g\circ f)(x))\in\mathrm{Gr}(g\circ f)$. Put $y=f(x)$, $z=g(y)$; then $(x,y,z)\in S$ and $p(x,y,z)=(x,z)=(x,(g\circ f)(x))$.

\sol{3.E.26(a)} Define $f:\Z\to\N$ by $f(n)=2n$ if $n\ge0$ and $f(n)=-2n-1$ if $n<0$ (values in $\N$: $2n\ge0$; for $n\le-1$, $-2n-1\ge1$). Inverse $g:\N\to\Z$: $g(m)=\frac m2$ if $m$ even, $g(m)=-\frac{m+1}2$ if $m$ odd. Check: $g(f(n))=g(2n)=n$ for $n\ge0$; for $n<0$, $f(n)=-2n-1$ is odd and $g(f(n))=-\frac{-2n-1+1}2=n$. And $f(g(m))=2\cdot\frac m2=m$ for even $m$; for odd $m$, $g(m)=-\frac{m+1}2<0$ and $f(g(m))=-2(-\frac{m+1}2)-1=m$. So $f$ is a bijection (Strategy 3.2.41).

\sol{3.E.26(c)} Let $A=\{0,1\}\cup\{\frac1n\mid n\ge2\}\subseteq[0,1]$ and $A'=\{\frac1n\mid n\ge2\}\subseteq(0,1)$. Define $f:[0,1]\to(0,1)$ by $f(0)=\frac12$, $f(1)=\frac13$, $f(\frac1n)=\frac1{n+2}$ for $n\ge2$, and $f(x)=x$ for $x\notin A$ (then $0<x<1$). Define $g:(0,1)\to[0,1]$ by $g(\frac12)=0$, $g(\frac13)=1$, $g(\frac1n)=\frac1{n-2}$ for $n\ge4$, and $g(y)=y$ for $y\notin A'$. Then $g(f(x))=x$ and $f(g(y))=y$ in every case ($f$ maps $A$ onto $A'$ by the shift $0,1,\frac12,\frac13,\dots\mapsto\frac12,\frac13,\frac14,\frac15,\dots$, and is the identity off $A$, where $[0,1]\setminus A=(0,1)\setminus A'$). So $f$ is a bijection.

\sol{3.E.26(d)} Compose bijections (Exercise 3.2.20): $[a,b]\to[0,1]$, $x\mapsto\frac{x-a}{b-a}$ (inverse $t\mapsto a+(b-a)t$); then $[0,1]\to(0,1)$ from (c); then $(0,1)\to(c,d)$, $t\mapsto c+(d-c)t$ (inverse $y\mapsto\frac{y-c}{d-c}$).

\sol{3.E.27} Write $T(k)=\binom{k+1}2=\frac{k(k+1)}2$, so $f(a,b)=T(a+b)+b$ and $T(k+1)=T(k)+(k+1)$. \emph{Claim:} for every $n\in\N$ there is a unique $k\in\N$ with $T(k)\le n<T(k+1)$. Existence: the set $\{k\in\N\mid T(k+1)>n\}$ is non-empty ($T(n+1)\ge n+1>n$), so it has a least element $k$ (well-ordering); then $T(k)\le n$ (if $k=0$, $T(0)=0\le n$; if $k>0$, minimality gives $T(k)\le n$). Uniqueness: $T$ is strictly increasing, so if $k<k'$ then $T(k+1)\le T(k')$, and $n<T(k+1)\le T(k')\le n$ is impossible. \emph{Surjective:} given $n$, take this $k$, put $b=n-T(k)$ and $a=k-b$. Then $0\le b<T(k+1)-T(k)=k+1$, so $b\le k$ and $a\in\N$; and $f(a,b)=T(k)+b=n$. \emph{Injective:} if $f(a,b)=f(a',b')=n$, then with $k=a+b$ we have $T(k)\le n=T(k)+b\le T(k)+k<T(k+1)$, so $k$ is the unique number of the claim; likewise $a'+b'=k$. Then $b=n-T(k)=b'$ and $a=k-b=a'$.

\sol{3.E.28} Let $Y=\{x\in X\mid e(x)=x\}$. Define $f:X\to Y$ by $f(x)=e(x)$: well-defined since $e(e(x))=e(x)$ shows $e(x)\in Y$. Define $g:Y\to X$ by $g(y)=y$. Then $g(f(x))=e(x)$, so $g\circ f=e$; and for $y\in Y$, $f(g(y))=e(y)=y$, so $f\circ g=\id_Y$.

\sol{3.E.29} False: $(1,1)$ and $(1,-1)$ are both in $G$, so uniqueness fails at $x=1$ (Theorem 3.1.9).

\sol{3.E.30} True: for each $x\in\Z$ the unique $y\in\N$ with $y^2=x^2$ is $y=|x|$ (if $y^2=x^2$ then $y=\pm x$, and only one of $x,-x$ is in $\N$ unless $x=0$). So $G=\mathrm{Gr}(x\mapsto|x|)$.

\sol{3.E.31} False: the empty function $\varnothing\to\{0\}$ (Definition 3.1.13) has empty domain and non-empty codomain.

\sol{3.E.32} True: this is Exercise 3.E.3.

\sol{3.E.33} False: $f:\{0\}\to\{1\}$, $f(0)=1$, has image $\{1\}\nsubseteq\{0\}$. (The image is a subset of the \emph{codomain}.)

\sol{3.E.36} True. A left inverse $g$ of $f$ is surjective by Exercise 3.2.28. If $g$ is a right inverse of $f$ ($f\circ g=\id_Y$) then $f$ is a left inverse of $g$, so $g$ is injective by Exercise 3.2.25.

\sol{3.E.39} Always: $f(x)=0$ for all $x\in X$ is well-defined, and any function $X\to\{0\}$ must take the value $0$ everywhere, so by function extensionality it equals $f$.

\sol{3.E.40} Sometimes. If $X=\varnothing$, the empty function $\varnothing\to\varnothing$ is the unique function (3.E.3). If $X\ne\varnothing$ there is no function $X\to\varnothing$ at all (3.E.3), so ``a unique function'' is false.

\sol{3.E.42} Never. The three values $f(1),f(2),f(3)$ lie in $\{1,2\}$, so two of them coincide: if $f(1)=f(2)$ then $(f(1),1),(f(1),2)\in G'$; otherwise $\{f(1),f(2)\}=\{1,2\}$ and $f(3)=f(k)$ for some $k\in\{1,2\}$, giving $(f(3),3),(f(3),k)\in G'$ with $k\ne3$. In all cases some first coordinate occurs twice in $G'$, so $G'$ is not a graph (Theorem 3.1.9).

\sol{3.E.43} Always: pick $x\in U$; then $f(x)\in f[U]$.

\sol{3.E.44} Sometimes. $f:\R\to\R$, $f(x)=x^2$: $V=\{-1\}$ gives $f^{-1}[V]=\varnothing$; $V=\{1\}$ gives $f^{-1}[V]=\{-1,1\}\ne\varnothing$.

\sol{3.E.45} Sometimes: $f(1)=1$, $f(2)=2$ is injective; $f(1)=f(2)=1$ is not.

\sol{3.E.46} Never: $f[\{1,2\}]=\{f(1),f(2)\}$ has at most two elements, so at least one of $1,2,3$ is not a value of $f$. (Precisely: if $f(1)=f(2)$ two of the three are missed; otherwise exactly the third element is missed.)
